% Generated by figures/13_14_24_25_granger_tables.R (score: a2) -- do not edit by hand.
% Requires \usepackage{booktabs} and \usepackage{threeparttable}.
\begin{table}[htbp]
\centering
\begin{threeparttable}
\caption{Granger causality between US and Chinese sentiment, by start date}
\label{tab:granger-time}
\small
\begin{tabular}{l rrr @{\hspace{1.5em}} rrr}
\toprule
 & \multicolumn{3}{c}{Full data} & \multicolumn{3}{c}{MFA / Presidential} \\
\cmidrule(lr){2-4}\cmidrule(lr){5-7}
Data after & US as & China as & Weeks & US as & China as & Weeks \\
 & response & response &  & response & response &  \\
\midrule
2002 (full series) & \textbf{$<$0.001} & \textbf{0.014} & 879 & \textbf{$<$0.001} & 0.128 & 648 \\
2003 & \textbf{$<$0.001} & \textbf{0.016} & 872 & \textbf{$<$0.001} & 0.125 & 644 \\
2004 & \textbf{$<$0.001} & \textbf{0.016} & 835 & \textbf{$<$0.001} & 0.102 & 612 \\
2005 & \textbf{$<$0.001} & \textbf{0.006} & 797 & \textbf{$<$0.001} & 0.051 & 580 \\
2006 & \textbf{$<$0.001} & \textbf{0.008} & 757 & \textbf{$<$0.001} & 0.075 & 545 \\
2007 & \textbf{$<$0.001} & \textbf{0.041} & 715 & \textbf{0.002} & 0.209 & 509 \\
2008 & \textbf{$<$0.001} & 0.129 & 669 & \textbf{0.003} & 0.434 & 480 \\
2009 & \textbf{$<$0.001} & 0.164 & 619 & \textbf{0.006} & 0.543 & 450 \\
2010 & \textbf{0.001} & 0.184 & 570 & \textbf{0.003} & 0.649 & 421 \\
2011 & \textbf{0.002} & 0.311 & 521 & \textbf{0.005} & 0.929 & 388 \\
2012 & \textbf{0.003} & 0.242 & 471 & \textbf{$<$0.001} & 0.884 & 365 \\
2013 & \textbf{$<$0.001} & 0.159 & 427 & \textbf{$<$0.001} & 0.862 & 347 \\
2014 & \textbf{0.002} & 0.231 & 382 & \textbf{$<$0.001} & 0.970 & 321 \\
2015 & \textbf{$<$0.001} & 0.189 & 341 & \textbf{$<$0.001} & 0.871 & 293 \\
2016 & \textbf{0.004} & 0.076 & 295 & \textbf{0.005} & 0.274 & 257 \\
2017 & \textbf{0.011} & 0.104 & 249 & \textbf{0.021} & 0.352 & 225 \\
2018 & \textbf{0.005} & 0.151 & 202 & \textbf{0.001} & 0.093 & 187 \\
2019 & \textbf{0.017} & 0.160 & 152 & \textbf{0.009} & \textbf{0.026} & 142 \\
2020 & 0.348 & 0.765 & 100 & 0.454 & 0.473 & 93 \\
\bottomrule
\end{tabular}
\begin{tablenotes}[flushleft]
\footnotesize
\item \textit{Notes:} Wald test $p$-values of bivariate Granger causality tests with 3 lags of each series. ``US as response'' tests whether lagged Chinese sentiment improves the prediction of US sentiment beyond lagged US sentiment; ``China as response'' tests the reverse. Bold: $p < 0.05$. Sentiment is the mean hybrid score (approach 2) of the documents in each period. Each pair of series is first tested for a unit root (augmented Dickey--Fuller, constant and trend, 5\% level, full sample). If both series are stationary, the test is run in levels; if neither is, on first differences; if only one is, by the Toda--Yamamoto procedure (both series in levels with one additional, untested lag of each). All $p$-values are from $F$-tests. Each row uses the weeks after 1 January of the year shown; the MFA series begins in October 2002, so the first row covers the full overlap. Weeks are seven-day bins counted from 1 January 2000; a week enters when both series have at least one document. Full data: US, Presidential Library documents and a random sample of 2,000 Congressional documents; China, MFA press conference answers and People's Daily articles. MFA / Presidential: US, dated Presidential Library documents; China, dated MFA press conference answers. Every weekly series is stationary, so all tests are in levels.
\end{tablenotes}
\end{threeparttable}
\end{table}
